% \chapter{Some Additional Material}

% Test for first appendix file. Please note that if you are adding a table or figure that takes up a whole page, you must have some text underneath your appendix title before using the next page.

% \begin{table}[h]
% \caption[Aliquam mi nisi]{Aliquam mi nisi, tristique at rhoncus quis, consectetur non mi. Phasellus blandit quam ligula, a viverra lacus commodo at.}
% %\begin{center}
% \begin{tabularx}{\textwidth}{XXXX}\hline
% Some    & Data  & Goes  & Here\\\hline
% Some    & Data  & Goes  & Here\\
% Some    & Data  & Goes  & Here\\
% Some    & Data  & Goes  & Here\\\hline
% \end{tabularx}
% %\end{center}

% \end{table}


% \chapter{Secondary Appendix Content}

% {\bf Paragraph headings.} There is no official fourth level heading. Do not use the Paragraph heading feature in LaTeX, simply apply the bold characteristic to the first few words of a paragraph followed by a colon or period.

% Additionally, the rules regarding italics and bolding does not apply to appendices. You can do whatever you'd like here mostly.

% \chapter{Expert Opinions}

% Test for third appendix file.



\chapter{Supplementary Material}

This appendix provides additional details that support the main text, including extended experimental data, code listings, and further analysis.

\section{Experimental Data}
Detailed tables and figures from the grid search for hyperparameter optimization and the learnable gate reduction experiments are presented here.

\section{Code Listings}
Selected excerpts from the DiffLogic and MiniNets implementations are provided for reference.

% \chapter{Additional Analysis}
% This section includes further discussion of the limitations, alternative approaches, and future improvements considered during this research.
